% ==============================================================================
% chapters/ch03_soil_analysis.tex — บทที่ 3
% ==============================================================================
\chapter{การประเมินชนิดดินและธาตุอาหาร}
\label{ch:soil_analysis}

ดินเป็นรากฐานสำคัญที่สุดของผลผลิตทางการเกษตร การประเมินสมบัติของดินอย่างแม่นยำและรวดเร็วช่วยให้เกษตรกรสามารถบริหารจัดการปุ๋ย น้ำ และการปรับปรุงโครงสร้างดินได้อย่างมีประสิทธิภาพ ในบทนี้จะอธิบายระเบียบวิธีวิจัยและการประยุกต์ใช้ปัญญาประดิษฐ์ในการประเมินดิน 2 รูปแบบหลัก ได้แก่ การจำแนกประเภทดินจากภาพถ่ายกายภาพ (Soil Visual Classification) และการวิเคราะห์ธาตุอาหารดินผ่านสเปกโทรสโกปีและการประมวลผลข้อมูลตาราง (Spectroscopy \& Tabular Analysis)

\academicsection{การจำแนกชนิดเนื้อดินจากภาพถ่ายกายภาพ}
\label{sec:soil_types_cv}

การจำแนกเนื้อดินด้วยคอมพิวเตอร์วิทัศน์อาศัยการสกัดคุณลักษณะเด่นด้านเนื้อสัมผัส (Texture) การสะท้อนแสง และโทนสีของดิน ในระบบได้จัดเตรียมชุดข้อมูลดินจริง 5 ชนิดหลัก ได้แก่

\begin{enumerate}
    \item \textbf{ดินดำอินทรีย์ (Black Organic Soil):} ดินที่มีปริมาณอินทรียวัตถุสูง สีดำคล้ำ กักเก็บความชื้นและธาตุอาหารได้ดีเยี่ยม
    \item \textbf{ดินขี้เถ้าภูเขาไฟ (Cinder / Volcanic Soil):} ดินที่มีรูพรุนสูง การระบายน้ำดี อุดมไปด้วยแร่ธาตุเฉพาะตัว
    \item \textbf{ดินลูกรังและดินแดง (Laterite / Red Soil):} ดินที่พบได้มากในภาคตะวันออกของไทย มีสัดส่วนของเหล็กและอะลูมิเนียมออกไซด์สูง มักมีปัญหาการยึดตรึงฟอสฟอรัส
    \item \textbf{ดินพีตหรือดินพรุ (Peat Soil):} ดินอินทรีย์ที่มีความเป็นกรดจัดและการย่อยสลายยังไม่สมบูรณ์
    \item \textbf{ดินเหลืองหรือดินร่วนปนทราย (Yellow Soil):} ดินเนื้อหยาบ มีการระบายน้ำเร็ว และมีธาตุอาหารในระดับปานกลางถึงต่ำ
\end{enumerate}

ภาพถ่ายดินเหล่านี้ถูกปรับขนาดมาตรฐานเป็น $224 \times 224$ พิกเซล และผ่านกระบวนการปรับปรุงข้อมูลภาพ (Data Augmentation) เพื่อจำลองสภาพแสงในแปลงเพาะปลูกจริง

\academicsection{การวิเคราะห์ธาตุอาหารด้วยสเปกโทรสโกปีคลื่นแสงสะท้อน (OSSL)}
\label{sec:soil_spectroscopy}

วิธีมาตรฐานดั้งเดิมในการตรวจวิเคราะห์ธาตุอาหารดินคือ การนำตัวอย่างดินส่งเข้าห้องปฏิบัติการเคมีเปียก (Wet Chemistry) ซึ่งใช้เวลานาน มีต้นทุนสูง และสร้างสารเคมีตกค้าง ในปัจจุบันงานวิจัยระดับโลกได้เปลี่ยนผ่านสู่การใช้ \textbf{การวัดการสะท้อนแสงเชิงสเปกตรัม (Reflectance Spectroscopy)} ทั้งในช่วงคลื่นตามองเห็น-อินฟราเรดย่านใกล้ (Visible-Near Infrared: Vis-NIR, ความยาวคลื่น 350--2,500 nm) และช่วงคลื่นอินฟราเรดย่านกลาง (Mid-Infrared: MIR, ความยาวคลื่น 4,000--600 cm$^{-1}$)

ระบบได้เชื่อมต่อข้อมูลตัวอย่างจริงจากคลังข้อมูลเปิด \textbf{OSSL (Open Soil Spectral Library)} ซึ่งพัฒนาโดยองค์กร Soil Spectroscopy for the Global Good โดยในแต่ละตัวอย่างดินจะมีแถบคลื่นสเปกตรัมสะท้อนจำนวน 2,152 แถบคลื่น จับคู่กับผลการตรวจวัดทางเคมีในห้องปฏิบัติการ ดังแสดงในตารางที่ \ref{tab:ossl_properties}

\begin{table}[htbp]
\caption{พารามิเตอร์สมบัติทางเคมีและกายภาพของดินที่ตรวจวัดในฐานข้อมูล OSSL}
\label{tab:ossl_properties}
\centering
\fontsize{9}{12}\selectfont
\begin{tabularx}{\textwidth}{llcX}
\toprule
\textbf{สัญลักษณ์ตัวแปร} & \textbf{ชื่อพารามิเตอร์ดิน} & \textbf{หน่วยวัด} & \textbf{ความสำคัญทางการเกษตร} \\
\midrule
\texttt{n.tot} & ไนโตรเจนรวม (Total Nitrogen) & \% & การเจริญเติบโตของลำต้นและใบ \\
\texttt{p.ext\_mel3} & ฟอสฟอรัสสกัดได้ (Mehlich-3 P) & mg/kg & การพัฒนาระบบรากและการออกดอก \\
\texttt{k.ext\_mel3} & โพแทสเซียมสกัดได้ (Mehlich-3 K) & mg/kg & ความแข็งแรงของเนื้อไม้และคุณภาพผล \\
\texttt{ph.h2o} & ปฏิกิริยาดิน (Soil pH ในน้ำ 1:1) & --- & การควบคุมการละลายตัวของธาตุอาหาร \\
\texttt{oc.tot} & คาร์บอนอินทรีย์ (Organic Carbon) & \% & ตัวชี้วัดความอุดมสมบูรณ์และจุลินทรีย์ดิน \\
\texttt{cec} & ความจุแลกเปลี่ยนแคตไอออน (CEC) & cmol$_\text{c}$/kg & ความสามารถในการอุ้มธาตุอาหารปุ๋ย \\
\texttt{clay.tot} & สัดส่วนอนุภาคดินเหนียว (Clay) & \% & การอุ้มน้ำและการระบายอากาศของดิน \\
\texttt{sand.tot} & สัดส่วนอนุภาคทราย (Sand) & \% & ความโปร่งและการแทรกชอนของรากพืช \\
\bottomrule
\end{tabularx}
\end{table}

\academicsection{การสร้างแบบจำลองแนะนำพืชและสูตรปุ๋ยจากข้อมูลตาราง}
\label{sec:crop_recommendation}

นอกเหนือจากการจำแนกภาพและสเปกตรัม ระบบยังบรรจุชุดข้อมูลตารางจริง \texttt{Crop\_\allowbreak recommendation.csv} จำนวน 2,200 แถวข้อมูล ซึ่งประกอบด้วยฟีเจอร์สำคัญ 7 ประการ ได้แก่ ระดับไนโตรเจน (N), ฟอสฟอรัส (P), โพแทสเซียม (K), อุณหภูมิ, ความชื้นสัมพัทธ์, ค่าความเป็นกรด-ด่าง (pH) และปริมาณน้ำฝนเฉลี่ย

ผู้เรียนสามารถนำข้อมูลตารางนี้ไปฝึกฝนโมเดลจำแนกพืชที่เหมาะสม (Crop Recommendation Classifier) ด้วยอัลกอริทึมต้นไม้ตัดสินใจ (Decision Tree / Random Forest) หรือ XGBoost เพื่อสร้างระบบอัจฉริยะที่สามารถรับค่าจากเซนเซอร์ IoT ภาคสนาม แล้วคำนวณพืชที่ควรปลูกพร้อมสูตรการปรับปรุงดินได้อย่างแม่นยำ
